\chapter{绪论}

\section{研究背景}

优化问题广泛存在于路径规划、资源调度、参数寻优、模型训练和工程设计等场景中。许多实际问题具有非线性、多峰值、约束复杂或解析梯度难以获取等特点，传统确定性方法在求解效率、稳定性和适用范围上可能受到限制。启发式优化算法能够通过群体搜索、随机扰动和迭代更新等方式在复杂搜索空间中寻找较优解，是计算机类算法研究方向毕业设计中常见的选题类型。

算法研究类论文不仅需要说明算法流程是否能够运行，还应给出明确的问题定义、基准算法、改进动机、参数设置、实验数据和评价指标。论文写作中应重点阐明算法改进的依据、核心步骤、复杂度或收敛性讨论，以及与对比算法之间的实验差异。

\section{研究内容}

本文示例以改进粒子群优化算法为背景，展示算法研究类本科毕业设计正文中常见的内容组织方式。论文正文通常需要覆盖以下内容：

\begin{custenum}
  \item 建立待求解优化问题的数学模型，明确目标函数、约束条件和评价指标；
  \item 分析基准粒子群优化算法的基本流程、参数含义和可能存在的不足；
  \item 设计动态权重、局部扰动或种群更新等改进策略，并说明关键步骤；
  \item 设计对比实验，对收敛速度、求解精度、稳定性和运行时间进行分析。
\end{custenum}

\subsection{论文结构}

全文结构如下：第1章介绍课题背景、研究内容和论文组织结构；第2章说明问题建模、算法框架、改进策略和实验设计，并给出图表和公式排版示例；第3章总结全文工作并说明后续可改进方向。参考文献、附录和致谢按学校格式要求排版。
